\section{Specyfikacja problemu decyzyjnego}

\subsection{Zmienne decyzyjne}

Definiujemy następujące zmienne decyzyjne:
\begin{itemize}
  \item $x_{i,t}$ - liczba wyprodukowanych jednostek produktu $i$ w miesiącu $t$
  \item $s_{i,t}$ - liczba sprzedanych jednostek produktu $i$ w miesiącu $t$
  \item $inv_{i,t}$ - stan magazynowy produktu $i$ na koniec miesiąca $t$
  \item $over_{i,t}$ - zmienna binarna określająca czy sprzedaż produktu $i$ w miesiącu $t$ przekracza 80\% pojemności rynku
\end{itemize}

gdzie:
\begin{itemize}
  \item $i \in \{1,2,3,4\}$ - indeks produktu
  \item $t \in \{1,2,3\}$ - indeks miesiąca (1: styczeń, 2: luty, 3: marzec)
\end{itemize}

\subsection{Ograniczenia}

\subsubsection{Ograniczenia czasowe maszyn}

Dla każdego miesiąca $t \in \{1,2,3\}$:

\begin{enumerate}
  \item Szlifierki (4 sztuki):
  \begin{equation}
  0.4 \cdot x_{1,t} + 0.6 \cdot x_{2,t} \leq 4 \cdot 384 = 1536 \text{ godzin}
  \end{equation}

  \item Wiertarki pionowe (2 sztuki):
  \begin{equation}
  0.2 \cdot x_{1,t} + 0.1 \cdot x_{2,t} + 0.6 \cdot x_{4,t} \leq 2 \cdot 384 = 768 \text{ godzin}
  \end{equation}

  \item Wiertarki poziome (3 sztuki):
  \begin{equation}
  0.1 \cdot x_{1,t} + 0.7 \cdot x_{3,t} \leq 3 \cdot 384 = 1152 \text{ godzin}
  \end{equation}

  \item Frezarka (1 sztuka):
  \begin{equation}
  0.06 \cdot x_{1,t} + 0.04 \cdot x_{2,t} + 0.05 \cdot x_{4,t} \leq 1 \cdot 384 = 384 \text{ godzin}
  \end{equation}

  \item Tokarka (1 sztuka):
  \begin{equation}
  0.05 \cdot x_{2,t} + 0.02 \cdot x_{3,t} \leq 1 \cdot 384 = 384 \text{ godzin}
  \end{equation}
\end{enumerate}

\subsubsection{Bilanse magazynowe}

Dla każdego produktu $i \in \{1,2,3,4\}$ i miesiąca $t \in \{1,2,3\}$:
\begin{equation}
inv_{i,t} = inv_{i,t-1} + x_{i,t} - s_{i,t}
\end{equation}

Z warunkami początkowymi:
\begin{equation}
inv_{i,0} = 0 \text{ dla } i \in \{1,2,3,4\}
\end{equation}

I końcowymi:
\begin{equation}
inv_{i,3} = 50 \text{ dla } i \in \{1,2,3,4\}
\end{equation}

\subsubsection{Ograniczenia rynkowe}

Dla każdego produktu $i \in \{1,2,3,4\}$ i miesiąca $t \in \{1,2,3\}$:
\begin{equation}
s_{i,t} \leq M_{i,t}
\end{equation}

Gdzie $M_{i,t}$ to maksymalna liczba sztuk produktu $i$, którą może przyjąć rynek w miesiącu $t$ (zgodnie z tabelą z zadania).

\subsubsection{Ograniczenia pojemności magazynu}

Dla każdego produktu $i \in \{1,2,3,4\}$ i miesiąca $t \in \{1,2,3\}$:
\begin{equation}
inv_{i,t} \leq 200
\end{equation}

\subsubsection{Ograniczenia dotyczące obniżki dochodu}

Dla każdego produktu $i \in \{1,2,3,4\}$ i miesiąca $t \in \{1,2,3\}$:
\begin{align}
s_{i,t} &\geq 0.8 \cdot M_{i,t} - M \cdot (1 - over_{i,t}) \\
s_{i,t} &\leq 0.8 \cdot M_{i,t} + M \cdot over_{i,t}
\end{align}

Gdzie $M$ to duża liczba (tzw. big-M).

\subsubsection{Nieujemność zmiennych}

\begin{equation}
x_{i,t}, s_{i,t}, inv_{i,t} \geq 0 \text{ dla } i \in \{1,2,3,4\}, t \in \{1,2,3\}
\end{equation}

\subsection{Funkcje oceny}

\subsubsection{Oczekiwany zysk}

\begin{equation}
E[Zysk] = \sum_{t=1}^{3} \sum_{i=1}^{4} E[R_i] \cdot s_{i,t} \cdot (1 - 0.2 \cdot over_{i,t}) - \sum_{t=1}^{3} \sum_{i=1}^{4} 1 \cdot inv_{i,t}
\end{equation}

Gdzie $E[R_i]$ to oczekiwany dochód ze sprzedaży jednostki produktu $i$, który należy wyznaczyć z rozkładu t-Studenta ograniczonego do przedziału [5; 12] z parametrami $\mu$ i $\Sigma$.

\subsubsection{Średnia różnica Giniego (miara ryzyka)}

\begin{equation}
\Gamma(x) = \frac{1}{2} \sum_{t'=1}^{T} \sum_{t''=1}^{T} |r^{t'}(x) - r^{t''}(x)| \cdot p^{t'} \cdot p^{t''}
\end{equation}

Gdzie:
\begin{itemize}
  \item $r^t(x)$ - realizacja zysku w scenariuszu $t$ dla decyzji $x$
  \item $p^t$ - prawdopodobieństwo scenariusza $t$
  \item $T$ - liczba rozważanych scenariuszy
\end{itemize}

